\documentclass[11pt]{article}
\usepackage[margin=2.2cm]{geometry}
\usepackage{amsmath}
\usepackage[T1]{fontenc}
\usepackage{lmodern}
\usepackage{hyperref}
\title{Supplementary Methods\\Temporal limits of protein-sequence transfer for cancer dependency prediction under unseen similarity components}
\date{}
\begin{document}
\maketitle
\small

\section*{S1. Outcome-access boundary and data roles}

DepMap 24Q4 CRISPR gene-effect data were treated as post-discovery development
data. Results from that release and a deliberately non-scientific engineering
pilot had been inspected before protocol version 1.4 was locked; no 24Q4 result
is described as prospective confirmation. DepMap 26Q3 was the designated
same-platform time-out release. Before the canonical 26Q3 outcome file was
downloaded or parsed for this analysis, a server-side freeze bound the protocol,
analysis lock, validator, complete 24Q4 campaign, consensus predictions and the
prespecified organoid identifiers by SHA-256. An append-only access ledger then
recorded acquisition, retrieval, user-authorized source verification and
validation in chronological order. The source receipt records provenance and
does not independently prove source authenticity or absence of access outside
this project workflow.

The organoid source had been accessed previously and was therefore defined as a
secondary, non-outcome-naive analysis. No endpoint, threshold, seed, fold, gene
mapping or failed gate was changed after observing 26Q3 outcomes. A post-access
compatibility amendment accepted the official CSV's empty first header as its
row-index label; it changed neither source bytes nor numeric values and is
disclosed as non-prospective.

\section*{S2. 26Q3 time-out evaluation and processing correction}

The primary temporal cohort comprised ModelIDs present in 26Q3 and absent from
the hash-bound 24Q4 cell-model list. Gene columns were accepted only in the
format \texttt{SYMBOL (ENTREZ\_ID)} and mapped by exact symbol, without synonym
rescue. A gene required a frozen 24Q4 consensus prediction and finite 26Q3
scores in at least 80\% of eligible new models. The frozen primary, selected
baseline and global-reference predictions were evaluated without refitting or
recalibration. Temporal confirmation required positive method $R^2$,
$\Delta R^2\geq0.01$, and a paired component-bootstrap 95\% lower bound above
zero.

Because 26Q3 included a processing correction, the protocol also specified a
secondary shared-model sensitivity. For each eligible gene, the median 26Q3
minus 24Q4 effect among exact shared ModelIDs was added to every frozen 24Q4
prediction and to the frozen reference before scoring on the unchanged new-model
cohort. This correction uses 26Q3 outcomes; it cannot replace the uncorrected
primary analysis or alter its gates.

\section*{S3. Organoid crosswalk and reproducibility controls}

The 26Q3 model list was compared with 256 organoid source labels frozen before
26Q3 outcome access. Exact string equality was the registered first-stage rule.
If identifier namespaces differed and no frozen authoritative crosswalk resolved
them, overlap had to remain unknown rather than be declared zero. In that case,
the time-out and organoid analyses could not be presented as independent
replications. After temporal scoring, identifiers were reconciled using the
26Q3 \texttt{Model.csv} and the published organoid data-availability table
(Supplementary Table 2 of Herranz-Ors et al.). Each published sample ID and
Sanger alias was matched exactly against cell-line names, HCMI IDs and Sanger
IDs across the entire 26Q3 metadata catalogue. Each accepted link identified one
model in each source. Linked ModelIDs were assigned to the frozen release
cohorts, and the organoid study's individual IDs were used to count related
samples. Unmatched entries remained unresolved. This retrospective audit used
metadata alone.

Each job and aggregate was published immutably with a completion record written
last. The aggregate was independently replayed on the designated server and
bound the complete roster, input files, metrics and consensus arrays by SHA-256.
The manuscript's numeric macros and primary table were generated only after
verifying the terminal completion schemas, file sizes and metrics hashes. Passing
software tests or integrity checks was not treated as evidence of novelty,
biological mechanism or submission readiness.

\section*{S4. Diagnostic re-run and exploratory probes}

The original campaign archives were produced on a designated analysis host whose
per-job prediction storage is no longer reachable. To answer reviewer requests,
the frozen implementation was re-run on a second Linux GPU server. The diagnostic
lock is identical in all scientific fields to the primary lock except for the
registered host identity; all input hashes, seeds, folds, grids, solver settings,
endpoints and gates are unchanged, and the code differs only in path and host
handling. The re-run reproduced the locked headline development metrics to
within $8.7\times10^{-10}$ and the locked 26Q3 cohort and metrics to the same
numerical precision. Diagnostic outputs are labelled post hoc; they explain the
locked result and do not replace it or alter any gate.

Distance strata were defined on the consensus nearest-training cosine distance
across all genes. For each stratum, method and baseline $R^2$ were recomputed
relative to the consensus global reference, with paired 95\% intervals from
resampling whole sequence components. Learned coefficients $(a_m,b_m)$ were read
from each job's final constrained fit together with the SLSQP status, iterations
and constraint residuals. Identity-threshold sensitivity used the registered
secondary contract: five outer folds, split seed 42 and training seed 42 at
20\%, 40\% and 50\% identity, each analysed separately and never pooled with the
primary 30\% result.

Additional post hoc probes used the same five folds (split seed 42, training
seed 42) and outer-test labels only for scoring. Candidate KNN and ridge
hyperparameters were selected on outer-training genes by leave-one-out or
three-fold component cross-validation. A James--Stein-style scalar shrinkage of
the KNN prediction toward the outer-training global mean was fitted on
outer-training genes and clipped to $[0,1.5]$. A sign-relaxed distance stack
retained only $\sum_m a_m=1$ and $\sum_m b_m=0$ and was fitted on out-of-fold
base predictions. The t6 layer probe recomputed the residue-chunked embedding at
hidden layers 2, 4 and 6 with the same pooling and reran the five-fold campaign.

Lineage-specific means were computed for every Oncotree lineage with at least 30
models. Each lineage mean was centred on the pan-context gene mean, and a ridge
model on the frozen embedding was fitted on outer-training genes and evaluated on
outer-test genes under the same five folds. GO biological-process annotations
were taken from the Open Targets 26.09 target release; the 500 largest positive
and 500 largest negative residuals were tested against the annotated background
with one-sided Fisher tests and Benjamini--Hochberg correction.

\end{document}
